\uuid{4g7B}
\titre{ Continuité et dérivées partielles en un point }

\niveau{}
\module{ Fonctions de plusieurs variables }
\chapitre{}
\sousChapitre{}
\theme{}
\auteur{Grégoire Menet}
\datecreate{ 2026-03-22}
\organisation{}
\difficulte{}

\contenu{

\texte{ 
On considère la fonction $f : \mathbb{R}^2 \to \mathbb{R}$ définie par
\[
f(x,y)=
\begin{cases}
\dfrac{y^2}{\sqrt{x^2+y^2}} & \text{si } (x,y)\neq(0,0),\\[1ex]
0 & \text{si } (x,y)=(0,0).
\end{cases}
\]
}

\begin{enumerate}

\item   \question{Quel est le domaine de définition de $f$ ?}
\indication{}
\reponse{La fonction est définie sur tout $\mathbb{R}^2$. En effet, pour $(x,y)\neq(0,0)$, on a $x^2+y^2>0$, donc $\sqrt{x^2+y^2}$ est bien définie et non nulle. De plus, la valeur de $f(0,0)$ est donnée séparément. Ainsi,
\[
D_f=\mathbb{R}^2.
\]}

\item   \question{La fonction $f$ est-elle continue en $(0,0)$ ?}
\indication{}
\reponse{Oui. Pour $(x,y)\neq(0,0)$,
\[
|f(x,y)|=\frac{y^2}{\sqrt{x^2+y^2}}.
\]
Or on a $y^2\leq x^2+y^2$, donc
\[
|f(x,y)|\leq \frac{x^2+y^2}{\sqrt{x^2+y^2}}=\sqrt{x^2+y^2}.
\]
Lorsque $(x,y)\to(0,0)$, on a $\sqrt{x^2+y^2}\to 0$, donc par encadrement
\[
f(x,y)\to 0=f(0,0).
\]
Ainsi, $f$ est continue en $(0,0)$.}

\item   \question{Les dérivées partielles de $f$ en $(0,0)$ existent-elles ? Si oui les calculer.}
\indication{}
\reponse{On calcule les dérivées partielles par définition.

Pour la dérivée partielle par rapport à $x$ :
\[
\frac{\partial f}{\partial x}(0,0)
=
\lim_{h\to 0}\frac{f(h,0)-f(0,0)}{h}.
\]
Or pour $h\neq 0$,
\[
f(h,0)=\frac{0^2}{\sqrt{h^2+0^2}}=0.
\]
Donc
\[
\frac{f(h,0)-f(0,0)}{h}=\frac{0-0}{h}=0,
\]
et par suite
\[
\frac{\partial f}{\partial x}(0,0)=0.
\]

Pour la dérivée partielle par rapport à $y$ :
\[
\frac{\partial f}{\partial y}(0,0)
=
\lim_{h\to 0}\frac{f(0,h)-f(0,0)}{h}.
\]
Or pour $h\neq 0$,
\[
f(0,h)=\frac{h^2}{\sqrt{h^2}}=\frac{h^2}{|h|}=|h|.
\]
Donc
\[
\frac{f(0,h)-f(0,0)}{h}=\frac{|h|}{h}.
\]
Cette quantité vaut $1$ si $h>0$ et $-1$ si $h<0$, donc la limite n’existe pas.

Ainsi :
\[
\frac{\partial f}{\partial x}(0,0)=0
\qquad\text{et}\qquad
\frac{\partial f}{\partial y}(0,0)\ \text{n’existe pas}.
\]}

\end{enumerate}

}